\documentclass[12pt]{article}
\usepackage[utf8]{inputenc}
\usepackage[T1]{fontenc}
\usepackage{hyperref}
\title{Bar\'e Language Resources for K-12 Curriculum Development}
\author{Compiled by Claude Sonnet 5.5 (claude-sonnet-5-5) for Andr\'as Kornai}
\date{}
\begin{document}
\maketitle

\section{Introduction}
Bar\'e (ISO 639-3 bae) is an Arawakan language of the upper Rio Negro (Venezuela and Brazil). Sources found in search describe it as extinct: its last fluent speaker reportedly died in 1993, and a few people remembered something of it in 2001. This conflicts with the survey's description of about 1,000 speakers, so the status should be checked. Speakers shifted to Nheengatu, Spanish and Portuguese. Closely related varieties include Guinau and Mandawaca.

\section{Primary Educational Resources}
None found.

\section{Linguistic and Reference Materials}
\begin{itemize}
\item Wikipedia entry on the Bar\'e language. \url{https://en.wikipedia.org/wiki/Bar%C3%A9_language}
\item Wikipedia entry on the related Guinau language. \url{https://en.wikipedia.org/wiki/Guinau_language}
\item Bar\'e document in the University of Missouri DICE Arawakan collection. \url{https://dice.missouri.edu/assets/docs/awarak/Bare.pdf}
\item WALS entry for Bar\'e. \url{https://wals.info/languoid/lect/wals_code_bae}
\end{itemize}

\section{Religious and Community Resources}
\begin{itemize}
\item Wikipedia entry on the Bar\'e people. \url{https://en.wikipedia.org/wiki/Bar%C3%A9_people}
\item Orinoco.org tribe page for the Bar\'e. \url{https://orinoco.org/?tribe=bare}
\end{itemize}

\section{Digital and Audio Resources}
\begin{itemize}
\item Endangered Languages Project entry. \url{https://endangeredlanguages.com/lang/1864}
\end{itemize}

\section{Curriculum Development Context}
Little material exists and the language may have no fluent speakers. Any curriculum would be a revitalization effort based on archival sources, such as the DICE document, and on the Bar\'e community.

\end{document}